% !TEX root = main.tex
% Included as D.4 after step_learning_shared_setup; no independent appendix section.
% Original Theorem 4(i),(ii) remain separate from this reduced-family extension.
\subsection{Full-prefix positional routing}
\label{app:weak-seeded-full-prefix-routing}

All causal sites compete in this setting, and the current-letter relative bias
starts at zero. The result and proof use the following seed conditions; the
feature equations and metrics are specified below.



\paragraph{Setting.}\label{app:lr-setting}
Fix $n\ge2$, $p=1/4$, $0<\eta<1/2$ and $0<\varepsilon\le1$.
Use the explicitly tied family \eqref{eq:lr-gates}--\eqref{eq:lr-feature-param}
with $c_0=r_{u,0}=0$, $\alpha_{z,0}<0$, $v_0,\rho_0>0$, and
$\alpha_{s,0}>|\beta_0|+|\alpha_{z,0}v_0|$. The fixed query and centered value direction
are as in \cref{app:weak-routing-family}. QK consists only of trainable
positional projections of orthogonal position embeddings; every visible
causal site participates in softmax. All relative biases start at zero.

Assume a complexity bound $\mathsf P\ge2$ bounds $n$, the widths,
$\varepsilon^{-1}$, all initial factor magnitudes, and the Frobenius norms
of the projected QK matrices. Also require
 $\min\{|\alpha_{z,0}|,v_0,\rho_0,\alpha_{s,0}-|\beta_0|-|\alpha_{z,0}v_0|\}\ge \mathsf P^{-1}.$ 
The initial base scores must have moderate spread: for some $S_0\ge0$,
 $|B_{i_tj}(0)|\le S_0,\qquad e^{2S_0}\le \mathsf P.$ 
Base scores may be bounded perturbations of equal scores.

\begin{lrtheorem}[Finite-accuracy learned positional routing]
\label{thm:lr-conditional-bridge}
In the setting above, there exist positive constant mobilities
$\nu_F,\nu_Q,\nu_R$, with $\nu_F^{-1},\nu_Q^{-1},\nu_R$
polynomial in $\mathsf P,1/\eta$, such that joint population gradient flow first
reaches $A_t\ge2/3$ at every supervised position and then
\begin{equation}\label{eq:lr-finite-theorem}
 \LRstep(T)\le\eta,\qquad T\le\operatorname{poly}(\mathsf P,1/\eta).
\end{equation}
During completion, $A_t\ge8/13$. The mobilities are fixed at initialization.
\end{lrtheorem}



\subsubsection{A weakly planted feature without a planted route}
\label{app:weak-routing-family}

We now specify the symmetry-reduced family used by the remaining results.
Fix the query letter. Let $s=(-1)^{s_{t-1}}$ and let attention depend on positions, not the
random letter identities.  The full-prefix pooled, centered value is
\begin{equation}\label{eq:lr-centered-value}
 Z_t=\sum_{j\in\mathcal J}\pi_{tj}(b_j-p),\qquad
 \pi_{tj}=\frac{\exp(S_{i_tj})}{\sum_{k\le i_t}\exp(S_{i_tk})},
 \quad \mathcal J=\{j_1,\ldots,j_n\}.
\end{equation}
Non-letter positions have zero \emph{value}, not zero attention mass.
A fixed token/position-aware linear value direction realizes the centered contributions
$b_j-p$ for prompt letters and zero contributions on the remaining visible
sites.  The trainable value parameter is its amplitude $v$. For the fixed query, assign its token the usual match
value $1-p$ and subtract $1-p$ using its position embedding; assign count
and separator tokens value zero. Prompt-letter positions have zero positional
correction. Orthogonal token and position registers realize this exactly.

Let $c,r_u$ be the prior/feature readouts, $\rho$ the paired output coefficient,
and $\alpha_s,\alpha_z$ the state/letter input coefficients. The ReLUs have preactivations
\begin{equation}\label{eq:lr-gates}
 r_+=\operatorname{ReLU}(\alpha_s s+\alpha_z vZ_t+\beta),\qquad
 r_-=\operatorname{ReLU}(-\alpha_s s+\alpha_z vZ_t+\beta).
\end{equation}
The binary logit (positive for state sign $+1$) is
\begin{equation}\label{eq:lr-factor-logit}
 f_t=cs+r_u\rho(r_+-r_-).
\end{equation}
Here $\beta$ is the \emph{MLP bias}, distinct from the attention
relative biases $\beta_\ell^{\mathrm{rel}}$ introduced below.
Every displayed previous state is available directly in the residual stream.

If
\begin{equation}\label{eq:lr-gate-condition}
 \alpha_s>|\beta|+|\alpha_z v|,
\end{equation}
then for all inputs with $|Z_t|\le1$,
\begin{equation}\label{eq:lr-feature-param}
 \begin{aligned}
 r_+-r_-&=s(d+{\alpha_z v} Z_t),& d&=\alpha_s+\beta,\\
 f_t&=s(C+gZ_t),& C&=c+r_u\rho d,& g&=r_u\rho {\alpha_z v}.
 \end{aligned}
\end{equation}

When $s=1$, only $r_+$ is active; when $s=-1$, only $r_-$ is active.
The inactive preactivation is strictly negative by
\eqref{eq:lr-gate-condition}.  Substitution gives the stated identity.

The two hidden units may be trained as an \emph{explicitly tied Euclidean
submanifold}: opposite count weights, equal centered-letter weights and
biases, and opposite output weights.  The induced Euclidean metric assigns
factor $1/2$ to a gradient update of a scalar appearing twice with equal
magnitude.  The mobilities below define flow in the displayed effective factor coordinates.
If a matrix realization uses a fixed direction of nonunit norm, its norm
factor must be absorbed into the assigned coordinate mobility.

\subsubsection{Exact cross-entropy routing signs}
\label{app:weak-routing-signs}

Let $\sigma(x)=(1+e^{-x})^{-1}$, $y_t'=1-b_t$ and
\begin{equation}\label{eq:lr-logistic-risk}
 \LRstep=\frac1{n-1}\sum_{t=2}^n\mathbb E
 \left[\log(1+e^{C+gZ_t})-y_t'(C+gZ_t)\right].
\end{equation}
This is exactly the binary cross-entropy of
\eqref{eq:lr-feature-param} because multiplying the binary logit and
label by $s$ cancels the previous-state sign.
Let $A_t=\pi_{t,j_t}$ and $\bar A=\frac1{n-1}\sum_{t=2}^n A_t$.

\begin{lrlemma}[Innovation selects the current source]
\label{lem:lr-signs}
For position-only attention,
\begin{equation}\label{eq:lr-initial-gate-gradient}
 \left.\frac{\partial\LRstep}{\partial g}\right|_{g=0}
 =p(1-p)\bar A>0
\end{equation}
for every $C$. At any finite logits and $g<0$, let $L_t$ denote the single-time population risk before averaging over $t$. Then
\begin{equation}\label{eq:lr-feature-source-signs}
 \pd{L_t}{\pi_{t,j_t}}<0,\qquad \pd{L_t}{\pi_{tj}}\ge0
 \quad(j\in\mathcal J,\;j\ne j_t),\qquad
 \pd{L_t}{\pi_{tj}}=0\quad(j\notin\mathcal J).
\end{equation}
If $|C|+|g|\le B$, then
\begin{equation}\label{eq:lr-source-bounds}
 \pd{L_t}{\pi_{t,j_t}}\le-2p(1-p)|g|\sigma(-B),\qquad
 0\le \pd{L_t}{\pi_{tj}}\le\frac{p(1-p) g^2\pi_{tj}}4\quad(j\ne j_t\text{ a letter}).
\end{equation}
\end{lrlemma}
\begin{proof}
Independence gives
\[
 \mathbb E[(b_t-p)Z_t]=p(1-p)A_t,\qquad
 \left.\pd{\LRstep}{g}\right|_{g=0}
 =\frac{p(1-p)}{n-1}\sum_{t=2}^n A_t.
\]
For the current letter, condition on all other bits and put
$U=C+g\sum_{j\in\mathcal J\setminus\{j_t\}}\pi_{tj}(b_j-p)$. Then
\[
 \pd{L_t}{\pi_{t,j_t}}=gp(1-p)\mathbb E
 \left[1+\sigma(U+gA_t(1-p))-\sigma(U-gA_tp)\right]<0.
\]
If $|C|+|g|\le B$, the bracket is at least $2\sigma(-B)$.
For an incorrect letter $j$, put
$U=C+g\sum_{k\in\mathcal J\setminus\{j\}}\pi_{tk}(b_k-p)$.
The target is independent of $b_j$, so
\[
 \pd{L_t}{\pi_{tj}}=gp(1-p)\mathbb E
 \left[\sigma(U+g\pi_{tj}(1-p))-\sigma(U-g\pi_{tj}p)\right].
\]
Since $g<0$ and $0<\dd{\sigma(u)}{u}\le1/4$,
\[
 0\le\pd{L_t}{\pi_{tj}}\le\frac{p(1-p)g^2\pi_{tj}}4.
\]
For $j\notin\mathcal J$, the value is zero and $\pd{L_t}{\pi_{tj}}=0$.
\end{proof}

The derivative with respect to a pre-softmax score is
\begin{equation}
 \pd{\LRstep}{S_{i_tj}}=\frac{\pi_{tj}}{n-1}
 \left(\pd{L_t}{\pi_{tj}}-\sum_{k\le i_t}\pi_{tk}\pd{L_t}{\pi_{tk}}\right).
\end{equation}

\begin{lrlemma}[Correct-versus-incorrect score comparison]
\label{lem:lr-postrouting-gradient}
Whenever $g<0$ and $A_t\ge1/2$,
\begin{equation}\label{eq:lr-score-comparison}
 \pd{\LRstep}{S_{i_t,j_t}}<0,\qquad \pd{\LRstep}{S_{i_tj}}-\pd{\LRstep}{S_{i_t,j_t}}\ge0
 \quad\text{for every other visible site }j.
\end{equation}
Consequently, if the attention biases are shared by relative offset and
$i_t-j_t=n$, then for every $\ell\ne n$,
\begin{equation}\label{eq:lr-bias-gap-mono}
 \frac{d}{dt}
 (\beta_n^{\mathrm{rel}}-\beta_\ell^{\mathrm{rel}})\ge0
\end{equation}
under common positive relative-bias mobility.
\end{lrlemma}
\begin{proof}
By \cref{lem:lr-signs},
\[
 \pd{L_t}{\pi_{t,j_t}}<0,\quad
 \pd{L_t}{\pi_{tj}}\ge0\ (j\ne j_t),\quad
 \sum_{k\le i_t}\pi_{tk}\pd{L_t}{\pi_{tk}}
 \ge A_t\pd{L_t}{\pi_{t,j_t}}.
\]
The current-letter derivative is strictly below the weighted mean, so
$\pd{\LRstep}{S_{i_t,j_t}}<0$.
Since $\pi_{tj}\le1-A_t\le A_t$, the softmax derivative identity gives
\begin{align*}
 &(n-1)\left(\pd{\LRstep}{S_{i_tj}}-\pd{\LRstep}{S_{i_t,j_t}}\right)\\
 &=\pi_{tj}\pd{L_t}{\pi_{tj}}-A_t\pd{L_t}{\pi_{t,j_t}}
 +(A_t-\pi_{tj})\sum_{k\le i_t}\pi_{tk}\pd{L_t}{\pi_{tk}}\\
 &\ge\pi_{tj}\pd{L_t}{\pi_{tj}}
 +A_t(1-A_t+\pi_{tj})\left(-\pd{L_t}{\pi_{t,j_t}}\right)\ge0.
\end{align*}
Offset $n$ occurs in every supervised row; another offset occurs at most once
per row. Summing over rows, including the negative correct-score derivatives
in rows where $\ell$ is absent, yields
\[
 \pd{\LRstep}{\beta_n^{\mathrm{rel}}}
 \le\pd{\LRstep}{\beta_\ell^{\mathrm{rel}}},\qquad
 \ddt{(\beta_n^{\mathrm{rel}}-\beta_\ell^{\mathrm{rel}})}
 =-\nu_R\left(\pd{\LRstep}{\beta_n^{\mathrm{rel}}}
 -\pd{\LRstep}{\beta_\ell^{\mathrm{rel}}}\right)\ge0.
\]
\end{proof}

\subsubsection{Feature bootstrap with every parameter group active}
\label{app:weak-routing-bootstrap}

All training uses one \emph{fixed} system of positive mobilities:
$c,r_u$ have mobility one; tied MLP/value factors have mobility
$\nu_F>0$; positional QK have mobility $\nu_Q>0$; all relative biases
have common mobility $\nu_R>0$.
No parameter group is explicitly frozen or switched at a chosen time.

Let $d=\alpha_s+\beta>0$ and $\alpha_z v<0$.

While the gate regime holds, the tied Euclidean gradient equations include
\begin{align}
 \ddt{c}&=-\pd{\LRstep}{C},\label{eq:lr-factor-u}\\
 \ddt{r_u}&=-\rho\left(d\pd{\LRstep}{C}+\alpha_z v\pd{\LRstep}{g}\right),\\
 \ddt{\rho}&=-\frac{\nu_F r_u}{2}
 \left(d\pd{\LRstep}{C}+\alpha_z v\pd{\LRstep}{g}\right),\\
 \ddt{\alpha_s}=\ddt{\beta}&=-\frac{\nu_F r_u\rho}{2}\pd{\LRstep}{C},\\
 \ddt{\alpha_z}&=-\frac{\nu_F r_u\rho v}{2}\pd{\LRstep}{g},\qquad
 \ddt{v}=-\nu_F r_u\rho\alpha_z\pd{\LRstep}{g}.
 \label{eq:lr-factor-bv}
\end{align}
In particular,
\begin{equation}\label{eq:lr-factor-invariants}
 r_u^2-\frac{2}{\nu_F}\rho^2=-\frac{2}{\nu_F}\rho_0^2,
 \qquad
 v^2-2\alpha_z^2=v_0^2-2\alpha_{z,0}^2,
 \qquad \alpha_s-\beta=\alpha_{s,0}-\beta_0.
\end{equation}
The first invariant gives $\rho^2\ge\rho_0^2$.  Under the gate
condition, the independent prior readout satisfies
$c(t)\le c_\star=\log((1-p)/p)$: at $c=c_\star$, with $r_u\rho\ge0$,
its derivative points inward.  At $r_u=0$ the velocity points into $r_u>0$,
and $\alpha_z v$ retains its negative sign as long as the bounded gate regime
holds: on either boundary $\alpha_z=0$ or $v=0$, one has $\pd{\LRstep}{g}=p(1-p)\bar A>0$,
so the vector field points into $\alpha_z<0,v>0$ (the simultaneous zero is excluded
by uniqueness from the nonzero seed). Thus $g=r_u\rho {\alpha_z v}<0$ for positive times.

\begin{lrlemma}[Quantitative weak-feature bootstrap, with explicit exits]
\label{lem:lr-feature-bootstrap}
Suppose on an interval prior to reaching the target magnitude that
\begin{equation}\label{eq:lr-amplification-region}
 \bar A\ge A_{\min}>0,\quad
 d\le M_d,\quad
 0<W\le|{\alpha_z v}|\le W_{\max},\quad
 \rho\ge\rho_{\min}>0,
\end{equation}
and the gate condition \eqref{eq:lr-gate-condition} holds.  Put
\begin{equation}\label{eq:lr-gamma-bootstrap}
 H_F=M_d+W_{\max},\quad
 \gamma=p(1-p) A_{\min}\left(\frac W{H_F}\right)^2,
 \quad
 T_{\mathrm{boot}}=\frac4{3\rho_{\min}^2H_F^2}.
\end{equation}
Starting with $c=r_u=0$, the quantity $x=-g$ reaches $x\ge\gamma$
by $T_{\mathrm{boot}}$, unless one of the region conditions fails first.
\end{lrlemma}
\begin{proof}
At $c\le c_\star$, since
$0<d+{\alpha_z v} Z_t\le H_F$, logistic curvature implies
$\pd{\LRstep}{C}\le{r_u\rho} H_F/4$.  At $g=0$, $\pd{\LRstep}{g}=p(1-p)\bar A$ and the
same curvature bound gives $\pd{\LRstep}{g}\ge p(1-p) A_{\min}-x/4$.
For $x\le\gamma$,
\begin{align*}
 d\pd{\LRstep}{C}+\alpha_z v\pd{\LRstep}{g}
 &\le\frac{M_dH_F}{4W}x-Wp(1-p)A_{\min}+\frac{W_{\max}}4x\\
 &\le-\frac34Wp(1-p)A_{\min}.
\end{align*}
The last step uses
$M_d H_F+WW_{\max}\le H_F^2$.
From \eqref{eq:lr-factor-u} and
$\rho^2\ge\rho_{\min}^2$,
$\ddt{(r_u\rho)}\ge(3/4)\rho_{\min}^2Wp(1-p) A_{\min}$.
In this region $\pd{\LRstep}{g}>0$ and
$\ddt{(\alpha_z v)}=-\nu_F{r_u\rho}(v^2+\alpha_z^2/2)\pd{\LRstep}{g}\le0$.
Therefore
\[
 \ddt{x}\ge\frac34\rho_{\min}^2W^2p(1-p)A_{\min},\qquad
 \frac{\gamma}{(3/4)\rho_{\min}^2W^2p(1-p)A_{\min}}
 =\frac4{3\rho_{\min}^2H_F^2}=T_{\mathrm{boot}}.
\]
\end{proof}

The groupwise loss budget is exact in the induced Euclidean metrics:
for the tied feature coordinates $(\alpha_s,\beta,\alpha_z,\rho,v)$, the squared
metric weights are $(2,2,2,2,1)$. For projected positional rows $(Q,K)$,
the squared metric weight is $\varepsilon^{-2}$.
Indeed, with orthogonal positional embeddings $\varepsilon e_i$,
projected rows satisfy $\ddt{Q}=-\nu_Q\varepsilon^2\nabla_Q L$ and
$\ddt{K}=-\nu_Q\varepsilon^2\nabla_K L$. This is the only QK geometry
used in the continuation argument.
\begin{equation}\label{eq:lr-mobility-budget}
 -\ddt{\LRstep}=
 \LRnorm{\nabla_{c,r_u}\LRstep}^2
 +\nu_F\LRnorm{\nabla_F\LRstep}^2
 +\nu_Q\LRnorm{\nabla_{Q,K}\LRstep}^2
 +\nu_R\LRnorm{\nabla_{\beta^{\mathrm{rel}}}\LRstep}^2.
\end{equation}
Here group norms and gradients use those induced metrics. For each group
with mobility $\nu$, Cauchy--Schwarz implies
\begin{equation}\label{eq:lr-movement-budget}
 \LRnorm{\theta_G(T)-\theta_G(0)}
 \le\sqrt{\nu T\,[\LRstep(0)-\LRstep(T)]}
 \le\sqrt{\nu T\log2}.
\end{equation}
Thus sufficiently small but positive $\nu_F,\nu_Q$ preserve prescribed
strict feature-gate and base-score margins over any known finite horizon.


\subsubsection{Relative-position learning and post-routing convergence}
\label{app:weak-routing-convergence}

Use positional scores
\begin{equation}\label{eq:lr-relative-scores}
 S_{i_tj}=B_{i_tj}+\beta_{i_t-j}^{\mathrm{rel}},
 \qquad
 B_{i_tj}=\frac{q_{i_t}^\top k_j}{\sqrt{d_h}},
 \qquad \beta_\ell^{\mathrm{rel}}(0)=0.
\end{equation}
The correct offset $n$ is not privileged at initialization.

\begin{lrlemma}[Correct-offset learning in the small-feature region]
\label{lem:lr-bias-learning}
Assume all base scores satisfy $|B_{i_tj}|\le S_A$,
all relative biases start at zero, and
$0<x=-g\le A_{\min}/4$. If $A_t\ge A_{\min}>0$,
then the correct score derivative is negative and all other visible-site
score derivatives are nonnegative:
\begin{equation}\label{eq:lr-smallscore-sign}
 \pd{\LRstep}{S_{i_t,j_t}}<0,\qquad \pd{\LRstep}{S_{i_tj}}\ge0\ (j\ne j_t).
\end{equation}
Consequently, for the shared relative biases,
\begin{equation}\label{eq:lr-beta-routes}
 \ddt{\beta_n^{\mathrm{rel}}}>0,\qquad
 \ddt{\beta_\ell^{\mathrm{rel}}}\le0\quad(\ell\ne n).
\end{equation}
They imply the uniform floor
\begin{equation}\label{eq:lr-floor}
 A_t\ge A_{\min}:=\frac{e^{-2S_A}}m,
 \qquad m:=\max_t\#\{j:j\le i_t\}\le3n.
\end{equation}
If the initial bootstrap reaches $x=\gamma$ and
$\gamma\le A_{\min}/4$, there is a polynomially bounded fixed bias
mobility $\nu_R$ that drives every $A_t$ to $2/3$ before $x$ leaves
$[\gamma/2,A_{\min}/4]$, provided the feature and base-score
neighbourhoods hold over that interval.
\end{lrlemma}
\begin{proof}
By Lemma~\ref{lem:lr-signs}, the correct value derivative is negative;
incorrect-letter derivatives are nonnegative.  The correct derivative has the logit-independent bound
 $-\pd{L_t}{\pi_{t,j_t}}\ge p(1-p) x(1-xA_t/4)\ge p(1-p) x(1-x/4),$ 
because $\sigma$ is $1/4$-Lipschitz. Hence
 $\sum_{k\le i_t}\pi_{tk}\pd{L_t}{\pi_{tk}}\le-p(1-p) x A_t(1-x/4)+p(1-p) x^2/4<0$ 
when $x\le A_{\min}/4$ and $A_t\ge A_{\min}$.  Hence the softmax derivatives have the stated signs.
The bias sign conclusion follows from the fixed offset
\eqref{eq:lr-constant-lag}.
With $\beta_n^{\mathrm{rel}}\ge0$, all other biases $\le0$,
and $|B|\le S_A$, softmax normalization yields
$A_t\ge e^{-2S_A}/m$.

For the time estimate, whenever some row has $A_t\le2/3$, the exact
correct-score derivative satisfies
$-\pd{\LRstep}{S_{i_t,j_t}}\ge c_0\frac1{n-1}p(1-p) xA_{\min}$
with $c_0=1/4$ in this small-$x$ region: indeed
$-\pd{\LRstep}{S_{i_t,j_t}}\ge\frac1{n-1} A_t(1-A_t)(-\pd{L_t}{\pi_{t,j_t}})$.
Therefore
$\ddt{\beta_n^{\mathrm{rel}}}
\ge c_0\nu_Rp(1-p)\gamma A_{\min}/[2(n-1)]$
as long as $x\ge\gamma/2$ and routing is incomplete.
All rows achieve $A_t\ge2/3$ when
$\beta_n^{\mathrm{rel}}\ge
2S_A+\log(2(m-1))$.
Thus the additional time is at most
 $T_{\mathrm{route}} \le \frac{2(n-1)[2S_A+\log(2(m-1))]} {c_0\nu_Rp(1-p)\gamma A_{\min}}.$ 
On a bounded feature neighbourhood, $|\ddt{x}|\le M_x$ with $M_x$
independent of $\nu_R$ (attention probabilities stay in $[0,1]$).
Choose the fixed $\nu_R$ large enough that this time is less than
$\min\{1,\gamma/(4M_x),(A_{\min}/4-\gamma)/(2M_x)\}$.
For $p=1/4$, the choice $\gamma=p(1-p) A_{\min}(W/H_F)^2$
automatically gives $\gamma\le3A_{\min}/16<A_{\min}/4$.
\end{proof}

\begin{lrlemma}[Post-routing invariance under bounded QK-score drift]
\label{lem:lr-post-routing}
Suppose $g<0$, $A_t(t_0)\ge2/3$ for all supervised $t$, and the
base-score contrasts change by at most $\log(5/4)$ after $t_0$:
 $|[B_{i_t,j_t}-B_{i_tj}](t)-[B_{i_t,j_t}-B_{i_tj}](t_0)| \le\log(5/4).$ 
As long as $g<0$ and the centered-feature family holds, all later
correct-position attention weights obey
\begin{equation}\label{eq:lr-post-floor}
 A_t(t)\ge\frac8{13}>\frac35.
\end{equation}
\end{lrlemma}
\begin{proof}
As long as all $A_t\ge1/2$, the bias comparisons of
Lemma~\ref{lem:lr-postrouting-gradient} imply that each relative-bias
gap $\beta_n^{\mathrm{rel}}-\beta_\ell^{\mathrm{rel}}$ is
nondecreasing.  The assumed bound on the other score contrasts says
the odds of each competitor relative to the correct site can increase
by at most a factor $5/4$.  Since the sum of competitor odds is at most
$1/2$ at time $t_0$, it stays at most $5/8$, yielding
$A_t\ge1/(1+5/8)=8/13$.  A first-exit argument closes the condition
$A_t\ge1/2$.
\end{proof}

\begin{lrlemma}[Uniform separability and convex-readout descent]
\label{lem:lr-readout-convergence}
Suppose $A_t\ge3/5$ at every supervised position.  Then
\begin{equation}\label{eq:lr-bit-separation}
 (1-2b_t)\left(\frac12-p-Z_t\right)\ge\frac1{10}
\end{equation}
for every prompt-bit realization.  Suppose in addition that, on a
finite interval after $t_0$, the trainable feature factors remain in a
neighbourhood that admits a \emph{fixed} readout comparator
$\widehat r=(\widehat c,\widehat r_u)$ satisfying
\begin{equation}\label{eq:lr-comparator}
 \LRstep(\widehat r;\phi(t))\le\eta/2
 \quad\text{for all }t\text{ in the interval}.
\end{equation}
If the two readout coefficients $(c,r_u)$ continue to train at unit
mobility and every other group has nonnegative mobility, then
\begin{equation}\label{eq:lr-descent}
 \LRstep(t_0+T)\le\frac\eta2+
 \frac{\LRnorm{r(t_0)-\widehat r}^2}{2T}.
\end{equation}
In particular, $T=\LRnorm{r(t_0)-\widehat r}^2/\eta$ suffices for
$\LRstep\le\eta$.
\end{lrlemma}
\begin{proof}
For $b_t=0$, $Z_t\le1-p-A_t$; for $b_t=1$,
$Z_t\ge A_t-p$.  Both cases imply
\eqref{eq:lr-bit-separation}.
At the reference seed, let $d_0=\alpha_{s,0}+\beta_0$ and $\alpha_{z,0}v_0<0$. Take
\[
 \begin{aligned}
 \widehat r_u&=\frac{20\log(2/\eta)}{\rho_0|\alpha_{z,0}v_0|},\\
 \widehat c&=20\log(2/\eta)\left(\frac12-p-\frac{d_0}{|\alpha_{z,0}v_0|}\right).
 \end{aligned}
\] 
At the reference feature, its effective logit is
$20\log(2/\eta)(1/2-p-Z_t)$, which has signed margin at least $2\log(2/\eta)$.
Require the feature neighbourhood to satisfy
\begin{equation}\label{eq:lr-comparator-stability}
 |\rho d-\rho_0d_0|+|\rho {\alpha_z v}-\rho_0{\alpha_{z,0}v_0}|
 \le\frac{\rho_0|{\alpha_{z,0}v_0}|}{20}.
\end{equation}
Since $|Z_t|\le1$, the logit perturbation is at most $\log(2/\eta)$.
The signed margin remains at least $\log(2/\eta)$, so the comparator loss is
at most $e^{-\log(2/\eta)}=\eta/2$. The same comparator is fixed before training,
independently of the routing hitting time.
At fixed feature factors, cross-entropy is convex in $r=(c,r_u)$ \citep[Section~3.1.5]{boyd2004convex}. Its first-order inequality \citep[Section~3.1.3]{boyd2004convex} gives
$\frac12\frac{d}{dt}\LRnorm{r-\widehat r}^2
\le\LRstep(\widehat r;\phi(t))-\LRstep(r;\phi(t))$.
The full joint flow makes the loss nonincreasing.  Integrating gives
\eqref{eq:lr-descent}.
\end{proof}

\begin{proof}[Proof with mobilities fixed before training]
Choose feature and QK balls of radii $r_F,r_Q>0$ around their initial
coordinates. Their inverse radii can be polynomially bounded in $\mathsf P$.
For the feature ball require the gate condition, $\alpha_z<0$, $v,\rho>0$,
\eqref{eq:lr-comparator-stability}, and uniform bounds
 $0<W\le|{\alpha_z v}|\le W_{\max},\quad d\le M_d,\quad \rho_{\min}\le\rho\le\Omega.$ 
The quantities $W^{-1},\rho_{\min}^{-1},W_{\max},M_d,\Omega$ are
polynomial in $\mathsf P$. For the QK ball require $|B|\le S_A=S_0+1$ and
 $|(B_{i_t,j_t}-B_{i_tj})-(B_{i_t,j_t}(0)-B_{i_tj}(0))| \le\tfrac12\log(5/4).$ 
Polynomial maps on bounded seed balls are Lipschitz with polynomial
constants; the assumed inverse-polynomial gate and factor margins give
these radii. Both QK points at any two times in the ball therefore have
contrast differences at most $\log(5/4)$.

Set $A_{\min}=e^{-2S_A}/m$, $H_F=M_d+W_{\max}$,
$\gamma=p(1-p) A_{\min}(W/H_F)^2$ and $T_b=T_{\mathrm{boot}}$ from
\eqref{eq:lr-gamma-bootstrap}. These quantities and their relevant
inverses are polynomial in $\mathsf P$, and $\gamma\le3A_{\min}/16$.
Set $T_{\mathrm{pre}}=T_b+1$. Until that horizon, inside the feature
ball, the unit-mobility readout equations give
 $|\ddt{c}|\le1,\qquad |\ddt{r_u}|\le\Omega H_F, \qquad \|(c,r_u)\|\le T_{\mathrm{pre}}\sqrt{1+\Omega^2H_F^2}.$ 
These bounds use only $|\sigma-y'|\le1$ and $|Z_t|\le1$ and do not
depend on bias mobility. With $\nu_F\le1$, the factor equations also
give a polynomial bound $M_x\ge1$ on $|\ddt{x}|$ in this interval,
uniformly over every attention distribution. Choose
 $\tau=\min\{1,\gamma/(4M_x),(A_{\min}/4-\gamma)/(2M_x)\}, \quad K_R=2S_A+\log(2(m-1)),$ 
 $\nu_R=\max\left\{1, \frac{2(n-1)K_R}{c_0p(1-p)\gamma A_{\min}\tau}\right\}, \qquad c_0=1/4.$ 
This fixes the relative-bias mobility in advance and makes the routing
bound at most $\tau$. It is polynomially bounded.

Use the fixed reference comparator from the readout lemma and define
 $R_{\mathrm{cmp}}=T_{\mathrm{pre}}\sqrt{1+\Omega^2H_F^2}+\|\widehat r\|, \qquad T_{\mathrm{post}}=R_{\mathrm{cmp}}^2/\eta,\qquad T_*=T_{\mathrm{pre}}+T_{\mathrm{post}}.$ 
Now fix the other mobilities by
 $\nu_F=\min\{1,r_F^2/(4T_*\log2)\},\qquad \nu_Q=\min\{1,r_Q^2/(4T_*\log2)\}.$ 
Their inverses are polynomially bounded. For the tied feature metric,
all squared weights are at least one; for projected QK the weight is
$\varepsilon^{-2}\ge1$. Thus the movement budget bounds ordinary
coordinate movement by $r_F/2$ and $r_Q/2$ up to $T_*$.
A first-exit argument therefore keeps both groups strictly inside their
chosen balls for the whole argument. Finite-time existence follows from
the smooth loss and its gradient-flow energy budget: all mobilities are
finite positive constants, so finite-time parameter escape is excluded.

Inside these balls, $c\le c_\star$ and $r_u\ge0$ persist by the inward
boundary signs stated above; $\rho>0$ and ${\alpha_z v}<0$ give $g<0$ for positive
times. In the small-feature region the bias signs imply the attention
floor by a first-exit argument from zero biases. The bootstrap lemma
reaches $x=\gamma$ by $T_b$ unless all routing is already complete;
even in that case the same bootstrap bound applies. Starting at this
hitting time, within at most $\tau$ the bias lemma makes every
$A_t\ge2/3$ (if this already holds, the extra time is zero). Let $t_0$
be the end of this interval; then $t_0\le T_{\mathrm{pre}}$.
The derivative bound keeps $x$ inside $[\gamma/2,A_{\min}/4]$ during
this routing interval. The bounded base-score contrasts and the
post-routing lemma preserve $A_t\ge8/13>3/5$ afterwards.
The feature ball preserves the fixed comparator for the entire remaining
interval. Its initial distance at $t_0$ is at most $R_{\mathrm{cmp}}$, so the readout
lemma gives $\LRstep(t_0+T_{\mathrm{post}})\le\eta$.
All mobilities were selected from initial uniform ball bounds; none
requires a later switch or a posterior choice of $t_0$.
\end{proof}

